\section{Proofs for policy-relevant continuation errors}\label{supp:contrast52}
\begin{proof}[Proof of Theorem~\ref{M-thm:contrast52}]
Fix a pass and suppress its index. For any feasible $a,b$, subtraction of the two Bellman expressions gives
\[
 Q_t(J_{t+1}^\pi;x,a)-Q_t(J_{t+1}^\pi;x,b)
 =Q_t(h_{t+1};x,a)-Q_t(h_{t+1};x,b)
       +\beta_t(P_t^a-P_t^b)e_{t+1}(x).
\]
Thus the accepted proposal has nonpositive true advantage relative to the incumbent. At a rejected state the implemented action is exactly the incumbent, whose advantage is zero. Backward induction gives $J^{\pi^+}\le J^\pi$. This proves safety even when only the proposal--incumbent contrast is certified.

For accuracy, suppose first that the proposal is accepted. For every feasible $a$, its approximate greediness and the uniform pair certificate imply
\[
 Q_t(J_{t+1}^\pi;x,v)-Q_t(J_{t+1}^\pi;x,a)
 \le\alpha_t+\beta_t\chi_t.
\]
Taking the infimum over $a$ requires neither existence of a minimizer nor a new selection assumption. If the proposal is rejected, then
$\overline d_t+\beta_tR_t(x;v,\pi_t)>0$ and
$\overline d_t\le d_t+\zeta_t$. Hence
\[
 Q_t(h_{t+1};x,\pi_t)
 <Q_t(h_{t+1};x,v)+\zeta_t+\beta_tR_t(x;v,\pi_t).
\]
Comparison with each feasible $a$ and another use of the action certificate yield an excess over $T_tJ_{t+1}^\pi$ of at most
$\alpha_t+\zeta_t+2\beta_t\chi_t$.
This common upper bound covers acceptance as well. Since future costs weakly decrease,
\[
 J_t^{\pi^+}=\mathcal T_t^{\pi^+}J_{t+1}^{\pi^+}
 \le\mathcal T_t^{\pi^+}J_{t+1}^{\pi}
 \le T_tJ_{t+1}^{\pi}+\tilde\epsilon_t.
\]
Monotonicity and the $\beta_t$ supremum-norm Lipschitz property of $T_t$ give the asserted recursion. The terminal gap is zero, and repeated substitution along successive passes gives the stated finite-horizon unrolling. Every retained policy is feasible for the original optimum, so a zero upper bound implies equality with that optimum. A fixed information contract may preclude zero greedy-search error; the theorem does not assume it away.

For sharpness take one period, two actions leading to distinct deterministic terminal states, zero current cost, and $\beta=1$. For $w>0$ and $0<\delta<w$, let the terminal critic be $(w,\delta)$ and the terminal cost be $(2w,\delta)$. The incumbent selects the first state and the exact greedy proposal the second. The exact action-error contrast is $w$, so $\chi=w$. The gate rejects because the critic contrast plus its allowance is $\delta>0$. The retained policy's gap is $2w-\delta$, whereas the terminal optimality gap and both search and arithmetic errors are zero. Letting $\delta$ decrease to zero excludes every universal coefficient below two.
\end{proof}

\begin{proof}[Proof of Proposition~\ref{M-prop:coupled52}]
The terminal assertion follows from $e_T=r_T$. At an earlier date, the policy-evaluation identity is
$e_t(x)=r_t(x)+\beta_tP_t^{\pi_t(x)}e_{t+1}(x)$.
The coupling's two marginals allow subtraction at states $x,y$ inside one integral. The triangle inequality and the inductive pair bound give
\[
 |e_t(x)-e_t(y)|\le H_t(x,y)+\beta_t
       \int D_{t+1}(u,v)\,d\Lambda_t^\pi(x,y)(u,v).
\]
The signed range supplies the separate bound $w_t$, so the minimum is valid. Coupling two action kernels at the same state gives the identical argument without the residual term and establishes the action certificate. Replacing an integral by an upper enclosure can only enlarge the bound before its intersection with the separately valid scalar width. This argument uses the marginals, not independence or optimality of the coupling.
\end{proof}

\begin{proof}[Proof of Corollary~\ref{M-cor:null52}]
The action difference annihilates $n$. Each of the two conditional means of $u$ lies in its verified interval, hence their difference in absolute value is at most its width. Further,
$J^\pi-(h+n)=u$, so the existing signed-residual recursion applied to the corrected critic yields the required interval without using an unknown true value. Adding $n$ to the continuation changes every feasible action objective by the same state-dependent amount; all action contrasts and approximate-greedy inequalities are therefore unchanged. The conclusion applies separately at each date and pass for which the membership and residual certificates are supplied.
\end{proof}

\begin{proof}[Proof of Proposition~\ref{M-prop:investment-null52}]
In the original dynamics the coefficients of $a$ in the first two coordinates are $1/2$ and $1/4$. For each common innovation $z$, the coefficient in $F_1(x,a,z)-2F_2(x,a,z)$ is consequently zero. Every function of that difference has identical realized values for two feasible actions with the same $z$, and thus identical conditional means under the original uniform innovation law. On the cube $y_1-2y_2+1$ ranges from $-1$ to $2$; its positive part ranges from zero to two. The global terminal error width is exactly $2M$. Its action contrast is nevertheless zero. The corrected critic is the original terminal cost $g$, for which the evaluation error vanishes. Proposal contrasts and their symbolic, outward-enclosed implementation are unchanged. Keeping the same proposals, tie conventions and cellwise gate therefore gives exactly the same terminal policy for every amplitude, not merely a similar observed cost.
\end{proof}

\subsection{Finite-state reference and limitations}
The exact reference implementation constructs a common-uniform coupling by exhausting ordered probability masses. It verifies both marginals, propagates pairwise residual bounds backward, and forms every same-state action contrast. Its regression suite independently evaluates the true finite-state policy values to test the certificates; those true values are not used by the certificate constructor. Tests cover one through five full improvement passes, a nonconstant action-null error with amplitudes zero, one, sixteen and 4096, non-null errors, constant shifts, exact feasibility, acceptance and rejection, and the sharp two-action example. These finite-state programs validate executable identities and counterexamples; they do not discretize the investment benchmark or establish full-horizon empirical convergence in that continuous economy.
